\documentclass[a4paper,11pt]{article}
\usepackage{exam-style-341}
\examinehead{341 农业综合知识三 · 数据库分章练习}{第01章\quad 基本概念与数据模型\quad 参考答案}
\begin{document}

\unitanswerscover{341 农业综合知识三 · 数据库分章练习}{第01章\quad 基本概念与数据模型 参考答案}{答案与解析按练习册题号顺序编排；判断题先给出判断，再给出改正后的正确表述；综合题给出完整的 E-R 要素、转换过程与两级映像分析。}

\answerpart{一、单项选择题（10 题，每题 2 分，共 20 分）}

\q \textbf{B}。\begin{analysis}数据库系统由数据库、数据库管理系统（DBMS）、应用程序、数据库管理员和硬件等组成，其中 DBMS 是位于用户与操作系统之间、用于科学地组织和存储数据、高效地获取和维护数据的系统软件，是整个数据库系统的核心。A 数据库只是被管理的数据集合；C 应用程序是用户访问数据库的手段；D 数据库管理员是数据库系统的人员组成。\end{analysis}

\q \textbf{B}。\begin{analysis}三级模式中，模式（逻辑模式）是数据库中全体数据的全局逻辑结构和特征的描述，一个数据库只有一个模式；外模式是用户能够看见和使用的局部数据的逻辑结构，一个数据库可以有多个；内模式是数据物理结构和存储方式的描述，一个数据库只有一个。D 存储模式是内模式的别名，与 C 同义。\end{analysis}

\q \textbf{A}。\begin{analysis}外模式/模式映像定义外模式与模式之间的对应关系。当模式改变时，DBA 只要修改该映像，就可以使外模式保持不变；由于应用程序是依据外模式编写的，因此应用程序不必修改，从而保证了数据的逻辑独立性。模式/内模式映像保证的是物理独立性。\end{analysis}

\q \textbf{C}。\begin{analysis}一个数据库只有一个模式、一个内模式，但可以有多个外模式：不同的用户或应用可以各自看到不同的局部数据视图，这正是外模式“面向用户”的体现。\end{analysis}

\q \textbf{A}。\begin{analysis}数据模型的三要素是数据结构、数据操作和完整性约束。B 是 E-R 图（概念模型）的三要素；C 是数据库的三级模式；D 是数据库系统的组成。注意“数据模型三要素”与“E-R 图三要素”不要混淆。\end{analysis}

\q \textbf{A}。\begin{analysis}概念模型是按用户的观点对数据和信息建模的模型，是现实世界到机器世界的中间层次，最常用的表示方法是 E-R 图。逻辑模型（如关系模型）和物理模型是概念模型向机器世界转换后的结果。\end{analysis}

\q \textbf{C}。\begin{analysis}数据管理技术经历了人工管理、文件系统和数据库系统三个发展阶段。人工管理阶段数据不保存、无共享、冗余大；文件系统阶段数据可以长期保存，但共享性差、冗余大、独立性差；数据库系统阶段由 DBMS 统一管理，数据结构化、共享性高、冗余度低、数据独立性高。\end{analysis}

\q \textbf{B}。\begin{analysis}非过程语言只说明“做什么”，不必指出“怎么做”，使用方便、可移植性好；SQL 是关系数据库的标准语言，属于非过程化语言。C、D 是低级语言。过程语言的代表是 C、Pascal 等。\end{analysis}

\q \textbf{D}。\begin{analysis}DBMS 的主要功能包括：数据定义（DDL）、数据操纵（DML）、数据库的运行管理（并发控制、安全性检查、完整性检查、恢复等）、数据库的建立和维护。编译并运行 C 源程序是编译系统的功能，不属于 DBMS。\end{analysis}

\q \textbf{B}。\begin{analysis}现实世界中的事物及其联系经过人的认识抽象为信息世界，再由信息世界转换为机器世界。概念模型（E-R 图）描述的是信息世界，是按用户观点对数据建模的中间层次；关系模型等逻辑模型描述机器世界。\end{analysis}

\answerpart{二、多项选择题（5 题，每题 4 分，共 20 分）}

\q \textbf{ABCD}。\begin{analysis}数据库系统一般由数据库、数据库管理系统（及其开发工具）、应用程序、数据库管理员和计算机硬件等组成。四个选项都是数据库系统的组成部分。\end{analysis}

\q \textbf{ABC}。\begin{analysis}数据库系统的主要特点是：数据结构化；数据共享性高、冗余度低、易扩充；数据独立性高；数据由 DBMS 统一管理和控制（安全性、完整性、并发控制和恢复）。D 错误：数据库系统中的数据由 DBMS 统一管理，不再由应用程序各自管理，这正是数据库系统区别于文件系统的重要标志。\end{analysis}

\q \textbf{ABCD}。\begin{analysis}三级模式中，模式是全体数据的全局逻辑结构，一个数据库只有一个；内模式是物理存储结构的描述，一个数据库只有一个；外模式是用户可见的局部逻辑结构，一个数据库可以有多个。四项均正确。\end{analysis}

\q \textbf{ABC}。\begin{analysis}数据模型由数据结构、数据操作和完整性约束三要素组成：数据结构描述系统的静态特性；数据操作描述系统的动态特性；完整性约束是一组完整性规则。D 数据独立性是数据库系统的特点，不属于数据模型的组成要素。\end{analysis}

\q \textbf{ABC}。\begin{analysis}E-R 图中实体用矩形表示，属性用椭圆表示，联系用菱形表示，实体与属性、实体与联系之间用无向边连接。D 错误：两个实体之间的联系类型有 1:1、1:n 和 m:n 三种，此外还存在三个或更多实体之间的多元联系。\end{analysis}

\answerpart{三、判断题（5 题，每题 2 分，共 10 分）}

\q √。\begin{analysis}数据库是长期存储在计算机内、有组织、可共享的大量数据的集合。数据库中的数据按一定的数据模型组织、描述和存储，具有较小的冗余度、较高的数据独立性和易扩展性。\end{analysis}

\q ×。改正：数据库系统的核心是数据库管理系统（DBMS），不是数据库。\begin{analysis}数据库只是被管理的数据集合；DBMS 是位于用户与操作系统之间、用于科学地组织和存储数据、高效地获取和维护数据的系统软件，数据库系统的一切操作都要通过 DBMS 完成，因此 DBMS 是数据库系统的核心。\end{analysis}

\q ×。改正：一个数据库只能有一个模式和一个内模式，但可以有多个外模式。\begin{analysis}模式是全体数据的全局逻辑结构，内模式是物理存储结构的描述，二者在一个数据库中都是唯一的；外模式面向不同用户或应用，可以有多个。\end{analysis}

\q ×。改正：模式/内模式映像保证数据的物理独立性；外模式/模式映像才保证数据的逻辑独立性。\begin{analysis}逻辑独立性是指模式改变时外模式和应用程序不变，由外模式/模式映像保证；物理独立性是指内模式改变时模式和应用程序不变，由模式/内模式映像保证。本题把两个映像的作用颠倒了。\end{analysis}

\q ×。改正：在 E-R 图中，一个联系可以涉及两个或两个以上实体，也可以在同一实体内部建立自联系。\begin{analysis}联系表示实体之间的关联。最常见的是两个实体之间的联系（二元联系），但也存在三个或更多实体之间的多元联系（如“供应商—零件—仓库”的供应联系），以及同一实体集内部的自联系（如“职工”实体中“领导”联系）。\end{analysis}

\answerpart{四、填空题（5 题，每题 2 分，共 10 分）}

\q 第一空：\textbf{人工管理}；第二空：\textbf{文件系统}。\begin{analysis}数据管理技术经历了人工管理、文件系统和数据库系统三个发展阶段。人工管理阶段数据不独立、不共享；文件系统阶段数据可以文件形式长期保存，但共享性差、冗余大；数据库系统阶段由 DBMS 统一管理。\end{analysis}

\q 第一空：\textbf{数据操作}；第二空：\textbf{完整性约束}。\begin{analysis}数据模型三要素：数据结构描述数据库的组成对象及对象之间的联系，是对静态特性的描述；数据操作是对数据库中各种对象的实例允许执行的操作的集合，是对动态特性的描述；完整性约束是一组完整性规则，用以保证数据的正确、有效和相容。\end{analysis}

\q 第一空：\textbf{模式}；第二空：\textbf{内模式}；第三空：\textbf{模式/内模式映像}。\begin{analysis}三级模式为外模式、模式、内模式；两级映像为外模式/模式映像和模式/内模式映像，分别保证逻辑独立性和物理独立性。\end{analysis}

\q 第一空：\textbf{属性}；第二空：\textbf{m:n}（多对多）。\begin{analysis}E-R 图三要素是实体、属性和联系；两个实体之间的联系类型有 1:1（一对一）、1:n（一对多）和 m:n（多对多）三种。\end{analysis}

\q 第一空：\textbf{逻辑}；第二空：\textbf{物理}。\begin{analysis}数据独立性包括逻辑独立性和物理独立性：前者由外模式/模式映像保证，后者由模式/内模式映像保证。\end{analysis}

\answerpart{五、简答题（2 题，每题 8 分，共 16 分）}

\q 参考答案要点：

\textbf{（1）三级模式（4 分）}
（1）外模式：又称子模式或用户模式，是数据库用户能够看见和使用的局部数据的逻辑结构，是用户的数据视图，一个数据库可以有多个外模式。
（2）模式：又称逻辑模式，是数据库中全体数据的全局逻辑结构和特征的描述，一个数据库只有一个模式。
（3）内模式：又称存储模式，是数据物理结构和存储方式的描述，一个数据库只有一个内模式。

\textbf{（2）两级映像（2 分）}
（1）外模式/模式映像：定义外模式与模式之间的对应关系。
（2）模式/内模式映像：定义模式与内模式之间的对应关系。

\textbf{（3）两级映像与数据独立性（2 分）}
（1）当模式改变时，DBA 修改外模式/模式映像，使外模式保持不变；应用程序依据外模式编写，因而不必修改，这就是数据的逻辑独立性。
（2）当内模式改变时，DBA 修改模式/内模式映像，使模式保持不变，从而外模式和应用程序都不必修改，这就是数据的物理独立性。
\begin{analysis}评分要点：三级模式的定义与数量（各 1 分，共 3 分）、两级映像的名称与作用（1 分）、两种独立性的说明（各 1 分，共 2 分）、模式/映像/应用程序三者的传递关系（2 分）。常见错误是把两种映像与两种独立性对应颠倒，或误认为一个数据库可以有多个内模式。\end{analysis}

\q 参考答案要点：

\textbf{（1）数据模型三要素（3 分）}
（1）数据结构：描述数据库的组成对象以及对象之间的联系，是对系统静态特性的描述。
（2）数据操作：对数据库中各种对象（型）的实例允许执行的操作的集合，包括操作及有关的操作规则，是对系统动态特性的描述。
（3）完整性约束：一组完整性规则，用以保证数据的正确、有效和相容。

\textbf{（2）概念模型的作用（2 分）}
概念模型是按用户的观点对数据和信息建模的模型，是现实世界到机器世界的中间层次。它独立于具体的 DBMS，便于用户与设计人员交流，是数据库概念结构设计阶段的产物，也是进一步转换为逻辑模型的依据。

\textbf{（3）E-R 图三要素（3 分）}
（1）实体：客观存在并可相互区分的事物，用矩形表示。
（2）属性：实体所具有的某一特性，用椭圆表示。
（3）联系：实体之间的关联，用菱形表示；两个实体之间的联系有 1:1、1:n、m:n 三种类型。
\begin{analysis}评分要点：三要素名称与含义（3 分）、概念模型是中间层次且独立于 DBMS（2 分）、E-R 图三要素及图形符号（3 分）。若只背“实体、属性、联系”而不说明概念模型的作用，最多得 5 分。\end{analysis}

\answerpart{六、综合题（2 题，每题 12 分，共 24 分）}

\q 参考答案：

\textbf{（1）实体、属性与联系类型（4 分）}
\begin{center}
\begin{tabular}{@{}llll@{}}
\toprule
对象 & 类型 & 属性 & 联系类型与说明 \\
\midrule
农户 & 实体 & 农户编号（标识）、姓名、联系电话 & 与作物为 m:n（种植） \\
作物 & 实体 & 作物编号（标识）、作物名称、单价 & 与订单为 m:n（订购） \\
订单 & 实体 & 订单号（标识）、订单日期 & 与农户为 1:n（下达） \\
种植 & 联系 & 面积 & 农户—作物，m:n \\
订购 & 联系 & 数量 & 订单—作物，m:n \\
下达 & 联系 & 无 & 农户—订单，1:n \\
\bottomrule
\end{tabular}
\end{center}

\textbf{（2）E-R 图（4 分）}

E-R 图的要素如下（考试作答时按下表画成图形：实体画矩形、属性画椭圆、联系画菱形，并在连线上标出 1、n、m）：
\begin{itemize}[leftmargin=1.6em, itemsep=1pt]
  \item 矩形实体“农户”，椭圆属性“农户编号（下加横线表示标识属性）”“姓名”“联系电话”；
  \item 矩形实体“作物”，椭圆属性“作物编号（标识属性）”“作物名称”“单价”；
  \item 矩形实体“订单”，椭圆属性“订单号（标识属性）”“订单日期”；
  \item 菱形联系“种植”连接“农户”和“作物”，农户端标 n、作物端标 m，菱形上连椭圆属性“面积”；
  \item 菱形联系“订购”连接“订单”和“作物”，订单端标 m、作物端标 n，菱形上连椭圆属性“数量”；
  \item 菱形联系“下达”连接“农户”和“订单”，农户端标 1、订单端标 n。
\end{itemize}

\textbf{（3）转换为关系模式（4 分）}
\begin{center}
\small
\begin{tabular}{@{}p{5.6cm}p{2.9cm}p{7.3cm}@{}}
\toprule
关系模式 & 主码 & 外码 \\
\midrule
农户（\underline{农户编号}，姓名，联系电话） & 农户编号 & 无 \\
作物（\underline{作物编号}，作物名称，单价） & 作物编号 & 无 \\
订单（\underline{订单号}，订单日期，农户编号） & 订单号 & 农户编号参照农户（农户编号） \\
种植（\underline{农户编号}，\underline{作物编号}，面积） & （农户编号，作物编号） & 农户编号参照农户；作物编号参照作物 \\
订购（\underline{订单号}，\underline{作物编号}，数量） & （订单号，作物编号） & 订单号参照订单；作物编号参照作物 \\
\bottomrule
\end{tabular}
\end{center}

转换规则：
（1）每个实体转换为一个关系模式，实体的属性作为关系的属性，实体的标识属性作为关系的主码。
（2）1:n 联系不单独建立关系模式，而是把 1 端实体的主码作为外码加入到 n 端实体的关系模式中（如“订单”中加入“农户编号”）。
（3）m:n 联系必须单独转换为一个关系模式，其属性由两端实体的主码加上联系自身的属性组成，主码为两端实体主码的组合，两端主码分别作为外码参照对应实体（如“种植”“订购”）。
（4）1:1 联系可以合并到任意一端的关系模式中，把另一端的主码作为外码加入，并对该外码加唯一性约束。
\begin{analysis}评分要点：三个实体及属性正确（1 分）；三个联系及其类型判断正确（2 分），其中“种植”“订购”为 m:n、“下达”为 1:n 是关键；E-R 图要素与基数标注正确（4 分）；转换后的五个关系模式及主码、外码正确（3 分）；转换规则表述完整（2 分）。常见错误：把 1:n 联系也单独建表（虽然可行但不是最简转换），或 m:n 联系漏掉联系属性“面积”“数量”，或忘记给“订单”加“农户编号”外码。\end{analysis}

\q 参考答案：

\textbf{（1）第（1）件事：模式改变，逻辑独立性（4 分）}

新增“客户风险等级”数据项是在全局逻辑结构上做修改，即\textbf{模式}发生了改变。此时需要修改\textbf{外模式/模式映像}，把新增的数据项按需要映射到相关外模式中；若原有外模式不需要包含该数据项，则外模式可以完全不变。由于应用程序是依据外模式编写的，只要外模式不变，应用程序就不必修改。这体现了数据的\textbf{逻辑独立性}。

\textbf{（2）第（2）件事：内模式改变，物理独立性（4 分）}

把数据文件从机械硬盘迁移到固态硬盘、改变存储组织和存取路径，改变的是数据的物理存储结构，即\textbf{内模式}发生了改变。此时需要修改\textbf{模式/内模式映像}，重新建立模式与新的内模式之间的对应关系，使模式保持不变。模式不变，则外模式和应用程序都不必修改。这体现了数据的\textbf{物理独立性}。

\textbf{（3）第（3）件事：新增外模式（2 分）}

为信贷部、柜面部分别建立视图，是新增\textbf{外模式}（用户视图）。每个外模式都需要通过外模式/模式映像与模式建立对应关系；新增外模式不影响原有的模式和内模式，也不影响其他外模式和应用程序。这体现了外模式面向不同用户、支持数据共享与安全保护的优点。

\textbf{（4）三级模式与两级映像的关系（2 分）}

三级模式是外模式、模式、内模式：外模式面向用户，模式是全局逻辑结构，内模式是物理存储结构。两级映像是外模式/模式映像和模式/内模式映像，它们把三层模式连接起来：外模式/模式映像保证逻辑独立性，模式/内模式映像保证物理独立性。映像由 DBMS 维护，模式或内模式改变时，DBA 通过修改相应映像使上层保持不变，从而把数据的变化与应用程序隔离，这正是数据库系统数据独立性高的根本原因。
\begin{analysis}评分要点：每件事的“哪一级模式—改哪个映像—应用程序是否修改—体现哪种独立性”四要素齐全各得 4 分（第（3）件事 2 分）；最后归纳映像与独立性的对应关系 2 分。易错点：把“增加数据项”误答为内模式改变，或把“更换存储设备”误答为模式改变；另外要注意“外模式/模式映像”与“模式/内模式映像”的书写顺序不能颠倒。\end{analysis}

\end{document}
